% Generated draft: requires mathematical and source review.
\documentclass[11pt]{article}
\usepackage[margin=25mm]{geometry}
\usepackage{fontspec}
\setmainfont{Latin Modern Roman}
\usepackage{amsmath,amssymb,mathtools,unicode-math}
\setmathfont{Latin Modern Math}
\usepackage{enumitem,needspace,xurl,hyperref}
\setlength{\parindent}{0pt}
\setlength{\parskip}{5pt}
\setlength{\emergencystretch}{3em}
\newcommand{\N}{\mathbb{N}}
\newcommand{\Z}{\mathbb{Z}}
\newcommand{\Q}{\mathbb{Q}}
\newcommand{\R}{\mathbb{R}}
\newcommand{\C}{\mathbb{C}}
\newcommand{\F}{\mathbb{F}}
\newcommand{\A}{\mathbb{A}}
\newcommand{\PP}{\mathbb{P}}
\newcommand{\E}{\mathbb{E}}
\newcommand{\eps}{\varepsilon}
\newcommand{\dd}{\,\mathrm{d}}
\begin{document}
\section*{UnsolvedMath ID: 30002746}
\textbf{Source problem code:} catalogue entry
\subsection{Definitions and mathematical statement}
The supplied source asks: \emph{Are different classroom participation patterns correlated with different trajectories of students' mathematical growth?}

A schematic notation for the objects mentioned is as follows. Let \(S\) denote a population of students, \(T\) an observation-time set, \(\mathcal P\) a space of classroom participation patterns, and \(M\) a space of mathematical-growth measurements. Write
\[
P:S\longrightarrow\mathcal P,
\qquad
G:S\longrightarrow M^T,
\]
where \(P(s)\) represents the participation pattern of student \(s\), and \(G(s)(t)\) represents that student's mathematical-growth measurement at time \(t\). These are placeholders for operational definitions, not additional hypotheses supplied by the source. In particular, participation may require a collective classroom description that is not captured by a student-level assignment alone.

\textbf{Research question.} Determine whether different classroom participation patterns are correlated with different trajectories of students' mathematical growth.

\textbf{Unresolved formulation.} The source does not specify \(S\), \(T\), \(\mathcal P\), \(M\), the measurement procedures, or the meaning of correlation. It supplies neither a sampling or probability model nor a particular association statistic, comparison criterion, or quantification over students, classrooms, and educational settings. Thus the question cannot presently be stated as a uniquely determined mathematical proposition or statistical hypothesis. No particular definition of correlation, causal claim, or universal law is asserted here.
\subsection{Short English statement}
The question asks whether students' classroom participation patterns are associated with how their mathematical understanding develops over time. The source leaves the relevant measurements, population, time scale, and meaning of association unspecified, so a precise mathematical or statistical claim requires further definitions.
\subsection{Sources}
\begin{itemize}
\item Local registry: \texttt{data/unsolvedmath/problems.json}, permanent UnsolvedMath ID \(30002746\), source problem code \texttt{catalogue entry}, titled \emph{Classroom Participation and Mathematical Growth Trajectories}. The supplied record labels the statement \texttt{exact}, describes it as internally readable and complete and retained byte-for-byte, and labels its status \texttt{open}. These are supplied registry qualifications, not independent verification of a precise mathematical formulation or its status.
\item Supplied report citation: \emph{Mathematics in Undergraduate Study Programs: Challenges for Research and for the Dialogue between Mathematics and Didactics of Mathematics} (2015), Oberwolfach Reports, pp.\ 3103--3175, DOI \texttt{10.4171/owr/2014/56}. Original OWR report: \url{https://doi.org/10.4171/owr/2014/56}. The record describes Oberwolfach Reports as documenting mathematical workshops and their open-problem discussions. It gives proposed year 2014 and identifies no individual proposer.
\item Supplied literature source: Tabach et al., knowledge shifts in mathematics classrooms, \url{https://www.sciencedirect.com/science/article/abs/pii/S0732312313001053}. The record supplies no fuller bibliographic details for this reference.
\item Supplied literature assessment, checked \texttt{2026-08-21}: status \texttt{open}. The assessment states that studies can coordinate individual and collective classroom analyses, but that no general law establishing the proposed correlation between participation patterns and mathematical-growth trajectories has been proved. This is supplied dated curation derived from the OWR curation, not an independently verified literature assessment. The record expressly calls for independent verification of time-sensitive status and cited claims.
\end{itemize}
\end{document}
